\documentclass{article}
\usepackage{amsmath, amssymb, amsthm}
\usepackage{geometry}
\geometry{a4paper, margin=1in}
\usepackage{fancyhdr}
\usepackage{hyperref}
\usepackage[english, russian]{babel}
\usepackage[utf8]{inputenc}
\usepackage{listings}
\usepackage{xcolor}
\usepackage{tikz}

\pagestyle{fancy}
\fancyhf{}
\fancyhead[L]{Algorithms}
\fancyhead[R]{\thepage}

\title{Algorithms 1 \\ Homework 11}
\author{Nikita Zagainov}
\date{\selectlanguage{english}\today}

\lstset{
    language=Python,
    basicstyle=\ttfamily\small,
    keywordstyle=\color{blue},
    stringstyle=\color{red},
    commentstyle=\color{gray},
    backgroundcolor=\color{lightgray!20},
    frame=single,
    rulecolor=\color{black},
    numbers=left,
    numberstyle=\tiny\color{gray},
    stepnumber=1,
    numbersep=10pt,
    showspaces=false,
    showstringspaces=false,
    breaklines=true,
    breakatwhitespace=true,
    tabsize=4,
    captionpos=b
}

\begin{document}

\maketitle

\begin{center}
	\section{Кратчайшие пути. Разные задачи.}
\end{center}
\begin{enumerate}
	\item Для начала для каждой вершины графа определим
	      набор из $k$ пар рёбер различного цвета. Тогда
	      при входе в эту вершину через одно ребро, мы выходим
	      через другое ребро другого цвета. Тогда, начиная с произвольной
	      вершины, мы гарантированно сможем пройти эйлеров цикл, заканчивающийся
	      в этой же вершине. Если этот цикл не обходит все вершины, то
	      выбираем новую вершину, принадлежащую циклу, и повторяем процедуру,
	      причём новая выбранная вершина теперь будет началом эйлерова цикла.
	      Эйлеров цикл существует в любом графе, у которого степень каждой вершины
	      чётная. Тогда решение цикл будет существовать только в том
	      случае, если существует такое распределение на пары. Очевидно, что
	      оно существует, если для каждой вершины не больше $k$ рёбер
	      одного цвета. Если это выполняется, то мы можем их спарить
	      следующим образом: В список длины $k$ через одно свободное место
	      ставим сначала самые частые цвета, далее - вторые по частоте
	      и так далее. Это гарантирует, что у нас не будет пар одного цвета,
	      так как тот факт, что каждого цвета у нас не более $k$, так
	      что в списке длины $2k$ они не могут стоять рядом. Далее,
	      просто пройдём по всем рёбрам, используя алгоритм Хиерхольцера.

	      Алгоритм:
	      \begin{lstlisting}
def compute_pairs(
    node: object,
    graph: dict[object, list[tuple[object, object]]],
    pairs: dict[object, dict[object, object]],
):
    color_lists = {}
    for neighbor, color in graph[node]:
        color_lists.setdefault(color, []).append(neighbor)
    node_cycle = [None] * len(graph[node])
    ptr = 0
    while color_lists:
        most_frequent_color = max(color_lists, key=lambda x: len(color_lists[x]))
        if len(color_lists[most_frequent_color]) > len(graph[node]) // 2:
            raise ValueError("No euler cycle")
        for neighbor in color_lists[most_frequent_color]:
            node_cycle[ptr] = (neighbor, most_frequent_color)
            ptr = ptr + 2
            if ptr >= len(graph[node]):
                ptr = 1
        del color_lists[most_frequent_color]
    node_pairings = {}
    n = len(node_cycle)
    for i in range(0, n, 2):
        a = node_cycle[i]
        b = node_cycle[i + 1]
        node_pairings[a] = b
        node_pairings[b] = a
    pairs[node] = node_pairings


def euler_cycle(graph: dict[object, list[tuple[object, object]]]) -> list[object]:
    pairs = {}
    for node in graph:
        if node not in pairs:
            compute_pairs(node, graph, pairs)
    path = []

    def visit(node):
        print(f"we are in {node}, our path so far is {path}, we need to inspect the following pairs: {list(pairs[node].keys())}")
        while pairs[node]:
            new_node, color = next(iter(pairs[node].keys()))
            pairs[node].pop((new_node, color))
            pairs[new_node].pop((node, color))
            print(f"we are in node {node}, going to {new_node} through {color}")
            visit(new_node)
        path.append(node)
        print(f"we are in {node}, our path so far is {path}, we finished visiting our pairs")

    start_node = next(iter(graph))
    visit(start_node)
    if any(pairs[node] for node in pairs):
        raise ValueError("No euler cycle")
    return path[::-1]
    \end{lstlisting}
	      Сложность алгоритма --- $O(|E|) = O(|V| + |E|)$. Сначала
	      для каждой вершины строим список пар за линейное время,
	      сумма длин всех этих списков --- $O(|E|)$. Далее, проходим
	      алгоритмом Хиерхольцера, который работает за $O(|E|)$, так как
	      гарантированно обходит каждое ребро один раз.

	\item Так как все рёбра в графе имеют вес, ограниченный сверху $W$,
	      то максимальное расстояние между вершинами графа не превосходит
	      $W(|V| - 1)$. Тогда можно хранить $W(|V| - 1)$ множеств, где $i$-е множество
	      содержит вершины, расстояние до которых от данной равно $i$. Чтобы
	      рассчитать расстояние до каждой вершины, для начала инициализируем
	      первое множество с начальной вершиной (расстояние равно 0), далее
	      проходим по этим множествам, на каждом шаге для каждой из вершин множества
	      считаем расстояние до её соседних вершин, и соседние вершины добавляем
	      в соответствующее множество (если оно уже не лежит в множестве, меньшем
	      по порядку). Это гарантирует, что мы найдём кратчайшие пути до каждой
	      из вершин, так как порядок рассмотрения вершин в нём аналоен порядку
	      алгоритма Дейкстры: на каждой итерации мы вибираем вершину с
	      наименьшим расстоянием до неё от истока.

	      Алгоритм:
	      \begin{lstlisting}
def shortest_paths(
    node: object,
    graph: dict[object, list[tuple[object, int]]],
    w: int,
) -> dict[object, tuple[int, list[object]]]:
    nodes = set()
    for key in graph:
        nodes.add(key)
        for edge in graph[key]:
            nodes.add(edge[0])

    sets = [set() for _ in range(w * (len(nodes) - 1) + 1)]
    sets[0].add(node)

    shortest_paths = {node: (0, [node])}

    for i in range(len(sets)):
        for node in sets[i]:
            for neighbor, weight in graph[node]:
                if neighbor not in shortest_paths:
                    sets[i + weight].add(neighbor)
                    shortest_paths[neighbor] = (
                        i + weight,
                        shortest_paths[node][1] + [neighbor],
                    )
                elif i + weight < shortest_paths[neighbor][0]:
                    sets[shortest_paths[neighbor][0]].remove(neighbor)
                    sets[i + weight].add(neighbor)
                    shortest_paths[neighbor] = (
                        i + weight,
                        shortest_paths[node][1] + [neighbor],
                    )

    return shortest_paths
    \end{lstlisting}

	\item
	      \begin{itemize}
		      \item[1.] В оригинальном алгоритме Дейкстры мы инициализируем кучу
			      с единственной вершиной с расстоянием 0, и остальными с
			      расстоянием $\infty$. Вершины с $\infty$ меняются только
			      после релаксации, так что если изменить алгоритм релаксации так,
			      что он добавляет вершину в кучу, если её там не было, то мы можем
			      не инициализировать кучу этими вершинами, так как алгоритм
			      релаксации вставляет недостающие вершины в кучу.

		      \item[2.] Модифицированный алгоритм может принимать во внимание
			      некоторые вершины по несколько раз, так что в случае, когда в
			      результате релаксации алгоритм обнаруживает, что какая-то вершина
			      с посчитанным расстоянием понизила своё значение, она будет добавлена
			      в очередь снова, и расстояния до них будут пересчитаны. Если же
			      в графе есть негативные циклы, то алгоритм не остановится, и будет
			      бесконечно уменьшать их расстояния. Если же негативных циклов нет,
			      этот алгоритм найдёт минимальные расстояния до всех вершин.
			      Сложность модифицированного алгоритма Дейкстры при рёбрах отрицательного
			      веса --- $O(|V|^2 \log |V| + |E||V|)$, так как данный алгоритм должен
			      в худшем случае проверить самые длинные рёберные пути,
			      равные $O(|V|)$ для каждой из $O(|V|)$ вершин, а это значит,
			      что каждая вершина может добавляться в кучу $O(|V|)$ раз, а
			      рёбра могут быть проверены суммарно для всех вершин $O(|E|)$
			      раз.

		      \item[3.] Для того, чтобы обнаруживать отрицательные циклы, можно
			      воспользоваться идеей алгоритма Беллмана-Форда и останавливать алгоритм,
			      если какая-то вершина была релаксирована более $|V| - 1$ раз.
	      \end{itemize}

\end{enumerate}
\end{document}


